\begin{scanpage}{0235}{217}
Our purpose was not to calculate the current on the body's surface, the field in the near zone, or the integral scattering thickness. These questions are investigated in a number of other works based on the physical theory of diffraction which were already enumerated in \S\ 25. In them, in particular, the first terms of asymptotic expansions in powers of $\lambda/a$ were obtained for the integral thickness which characterizes the total power scattered by a body. However, in these works, as a rule, equations are missing for the scattering characteristics which are valid with any direction of irradiation and observation. Therefore, the results of this book and the indicated works mutually supplement one another.

At present, only a limited number of diffraction problems have yielded to theoretical studies, as a result of which experimental studies of diffraction by various bodies have taken on great importance. In Chapter VI an experimental method was discussed which enabled one to isolate in a ``pure form'', and to measure, the field from the nonuniform part of the current excited by a plane wave on a metal body of any shape. In the same chapter, it was shown that the well-known phenomenon of depolarization of the wave reflected from a body which is found in free space is produced by the nonuniform part of the current, or, in other words, by the surface distortion.

The investigation carried out in Chapter VII for the problem of diffraction by a thin, finite length cylindrical conductor represents a natural development and completion of the method of considering the multiple diffraction of edge waves which was applied in Chapter V. In Chapter VII equations were derived for the scattering diagram which are suitable for vibrators of an arbitrary length with any irradiation and observation directions.

The results obtained in this book show the fruitfulness of physical diffraction theory, and enable one to arrive at the solution of other more complicated problems. Such problems may be divided into two classes. Problems which may now already be solved on the basis of the known results of diffraction theory are related to the
\end{scanpage}
